\begin{lemma}[Lemma 9n: available colors along the full tail]
Let $(T,\varphi_n)$ be a positive-rung extended Tashkinov tree, let
$T_n$ be its last ladder tree, and let
$T_n=T_{n,0}\subset\cdots\subset T_{n,q+1}=T$ be a good hierarchy with
base $T_{n,0}^{*}$.  Put
\[
 A=T_{n,q}^{*},\qquad \ell=|T|-|A|,\qquad
 D_{n,q}=\noderef{HierarchyStoredColorsAt}(H,q),
 \qquad
 \Gamma^q=\bigcup_{\eta\in D_{n,q}}\Gamma^q_\eta.
\]
Here $A$ is \noderef{HierarchyStarredStage} at the last divider and
$\ell$ is the number of exact tail additions above $A$.  For a natural number $i$, let
$T(y_i)$ be the ordered prefix of $T$ ending after the first $i$ entries
above $A$; here and below $|A|$ denotes its tree-sequence length, and
formally the prefix length is
$\min\{|A|+i,|T|\}$.  Define
\[
 B=\bigcup_{v\in V(T_{n,0}^{*})\setminus V(T_n)}
       \overline\varphi_n(v),
 \qquad
 U_i=\{\varphi_n(e_r): |A|\le r<|A|+i\},
\]
where $U_i$ contains only colors actually assigned to the displayed exact
edge identities in that index interval, and put
\[
 X_i=\overline\varphi_n(T(y_i))\setminus B\setminus U_i.
\]
Assume the exact prefix-elementarity assertion (6.6), namely that
$V(T(y_{\ell-1}))$ is elementary with respect to $\varphi_n$, and assume
\[
 2n+11\le |\overline\varphi_n(T_n)|,
 \qquad n=\mathit{ett.rounds}.
\]
Then, for every $0\le i\le \ell-1$ (equivalently, every natural number
$i<\ell$),
\[
                         |X_i|\ge 2n+11.                 \tag{1}
\]
Furthermore, if
\[
             |X_i\setminus(\Gamma^q\cup D_{n,q})|\le4, \tag{2}
\]
then there is an injective family of colors
$(\eta_h)_{h\in\operatorname{Fin}7}$ such that, for every $h$,
\[
 \eta_h\in D_{n,q}\cap\overline\varphi_n(T(y_i))
 \quad\hbox{and}\quad
 (\Gamma^q_{\eta_h}\cup\{\eta_h\})\cap U_i=\varnothing.
\]
\end{lemma}
\begin{proof}
Write $m=|A|$.  We first verify that all displayed
prefixes and edge indices in the statement are the paper's actual ones.
The nested and terminal-stage fields of \noderef{TashkinovHierarchy}
show that $A$ is an ordered initial segment of $T$.  Indeed, if $q>0$,
the last nested inclusion is $T_{n,q}^{*}=T_{n,q}\subset T_{n,q+1}=T$;
if $q=0$, this is the core-inside-first field followed by the terminal-stage
identity.  Thus $m\le |T|$.

From $m\le |T|$ and the definition $\ell=|T|-m$, we have
$m+\ell=|T|$.  Now fix $i<\ell$.  This inequality gives
$\ell>0$ and $i\le \ell-1$, while
\[
              m+i<m+\ell=|T|.
\]
Consequently the minimum in the formal definition of $T(y_i)$ is
$m+i$, and the one in $T(y_{\ell-1})$ is $m+\ell-1$.  This proves rather than
assumes both the positivity and the prefix-bound side conditions.  By
\noderef{TreeSequencePrefix}, $T(y_i)$ is the entry-by-entry prefix of
$T(y_{\ell-1})$, so (6.6) and
\noderef{ColoredVertexSetElementary} imply that $V(T(y_i))$ is elementary.

We next expand the first count.  Let
\[
 N_i=V(T(y_i))\setminus V(A).
\]
The vertex and edge maps of a \noderef{TreeSequence} are injective.
Since the two prefixes have respective lengths $m+i$ and $m$, their
vertex-set difference has exactly $i$ members.  The exact tail-edge
identities are $T.\mathit{edge}(r)$ for $m\le r<m+i$; there are also
exactly $i$ of them, and hence their set of colors $U_i$ has cardinality
at most $i$.  This uses the precise half-open interval appearing in the
Lean statement, so neither the distinguished first edge nor the next edge
after $T(y_i)$ is counted.

Every vertex misses at least one color.  To see this from the stated ETT
data, the index-zero certificate in \noderef{TashkinovSeries} is the
initial constructor of \noderef{TashkinovSeriesStep}, hence supplies the
\noderef{KTriple} from which the series starts.  Its degree clause is
$k\ge\Delta(G)+1$.  For any vertex $v$, choose for each member of
\noderef{PresentColors} at $v$ an incident edge bearing that color.
Two different colors cannot choose the same edge, because an edge has one
value under the partial coloring.  Thus the number of present colors is at
most the number of incident exact edge identities, namely
$d_G(v)$ by \noderef{Degree}, and this is at most
\noderef{MaximumDegree}.  Since $d_G(v)<k$, the complement definition
\noderef{MissingColors} shows that $\overline\varphi_n(v)$ is nonempty.

The missing-color sets at distinct vertices of $T(y_i)$ are pairwise
disjoint by elementarity.  Therefore the union definition
\noderef{TreeMissingColors}, applied to the $i$ vertices of $N_i$, gives
\[
 \left|\bigcup_{v\in N_i}\overline\varphi_n(v)\right|
   =\sum_{v\in N_i}|\overline\varphi_n(v)|\ge i\ge |U_i|.       \tag{3}
\]
Since $A$ is an initial segment of $T(y_i)$, the same disjointness also
gives
\[
 |\overline\varphi_n(T(y_i))|
 =|\overline\varphi_n(A)|+
   \left|\bigcup_{v\in N_i}\overline\varphi_n(v)\right|.       \tag{4}
\]

It remains to compare the base with $T_n$.  By
\noderef{HierarchyBase}, if the last extension is not PE then
$T_{n,0}^{*}=T_n$ and $B$ is empty.  In the PE case its
\noderef{OrderedTreeUnion} witness makes $T_n$ an ordered initial segment
of $T_{n,0}^{*}$.  Hence in both cases
$V(T_n)\subseteq V(T_{n,0}^{*})\subseteq V(A)$.  Elementarity and the
definition of $B$ yield the disjoint equality
\[
 |\overline\varphi_n(T_{n,0}^{*})|
   =|\overline\varphi_n(T_n)|+|B|.
\]
As the base is contained in $A$, this implies
\[
                  |\overline\varphi_n(A)|-|B|
                       \ge |\overline\varphi_n(T_n)|.           \tag{5}
\]
For arbitrary finite sets $Y,B,U$, the inclusion
$Y\subseteq((Y\setminus B)\setminus U)\cup B\cup U$ gives
$|((Y\setminus B)\setminus U)|\ge |Y|-|B|-|U|$ whenever the
right side is read through the corresponding cardinal inequality.
Apply this with $Y=\overline\varphi_n(T(y_i))$, and use (3)--(5):
\[
\begin{aligned}
 |X_i|
 &\ge |\overline\varphi_n(T(y_i))|-|B|-|U_i|\\
 &\ge |\overline\varphi_n(A)|-|B|\\
 &\ge |\overline\varphi_n(T_n)|\\
 &\ge 2n+11.
\end{aligned}
\]
This proves (1).

Assume now (2), and set
\[
 C=X_i\cap(\Gamma^q\cup D_{n,q}),\qquad d=|D_{n,q}|.
\]
The disjoint decomposition of $X_i$ into $C$ and the set in (2), together
with (1), gives
\[
                              |C|\ge 2n+7.                       \tag{6}
\]
By \noderef{HierarchyStoredColorsCardinality}, $d\le n$.

For each $\eta\in D_{n,q}$ put
$Q_\eta=\Gamma^q_\eta\cup\{\eta\}$.  We check the block structure
directly from the good-hierarchy fields.  The `twoColors' field of
\noderef{GoodHierarchy} says $|\Gamma^q_\eta|=2$, and its
`disjointPairs' field says that the Gamma pairs for distinct stored colors
are disjoint.  Its `avoidance' field puts every color of
$\Gamma^q_\eta$ among the colors missing at the last relevant hierarchy
stage.  For $q>0$ that stage is $A$; for $q=0$ it is $T_n$, whose vertex
set is contained in $A$.  On the other hand, unfolding
\noderef{HierarchyStoredColorsAt} gives
\[
 D_{n,q}=D_n\setminus\overline\varphi_n(A),
\]
so every member of $D_{n,q}$ is not missing on $A$.  It follows both that
$\eta\notin\Gamma^q_\eta$ and that no stored color belongs to any other
Gamma pair.  Therefore the blocks $Q_\eta$ are pairwise disjoint and each
has exactly three members.  Moreover
\[
 \Gamma^q\cup D_{n,q}=\bigcup_{\eta\in D_{n,q}}Q_\eta,
 \qquad |\Gamma^q\cup D_{n,q}|=3d\le3n.                         \tag{7}
\]
Equations (6) and (7) first imply $2n+7\le3n$, hence $n\ge7$.
They also give $3d\ge |C|\ge2n+7\ge21$, so $d\ge7$; in particular
there really are at least six blocks available for the padded estimate
below.

Call a block full when $Q_\eta\subseteq C$.  Suppose at most six blocks
were full.  A full block contributes at most three colors to $C$, while a
nonfull three-element block contributes at most two.  The blocks are
pairwise disjoint and $C$ is contained in their union.  Since $d\le n$
and $n\ge7$, padding to six full blocks and $n-6$ remaining blocks gives
the paper's exact estimate
\[
             |C|\le 3\cdot6+2\cdot(n-6)=2n+6,
\]
contradicting (6).  Thus at least seven distinct stored colors index full
blocks.  Choose seven of them and index them injectively as
$(\eta_h)_{h\in\operatorname{Fin}7}$.

For every $h$, fullness gives
$Q_{\eta_h}\subseteq C\subseteq X_i$.  In particular
$\eta_h\in D_{n,q}$ and, because $X_i$ is a subset of
$\overline\varphi_n(T(y_i))$, the color $\eta_h$ is missing on
$T(y_i)$.  Finally, the definition
$X_i=(\overline\varphi_n(T(y_i))\setminus B)\setminus U_i$
shows that every member of the whole block
$\Gamma^q_{\eta_h}\cup\{\eta_h\}$ avoids $U_i$.  This is precisely the
two asserted properties of the seven-color family.
\end{proof}
